\documentclass[10pt]{beamer}
\usetheme[]{Feather}
\setbeamercolor{Feather}{fg=black!90,bg=teal}
\setbeamercolor{structure}{fg=black}
\usepackage[utf8]{inputenc}
\usepackage[spanish]{babel}
\usepackage[T1]{fontenc}
\usepackage{helvet}
\usepackage{lipsum}
\usepackage{amsmath,amssymb,amsfonts}
\usepackage{amsthm}
\usepackage{physics}
\usepackage{framed, color}
\usepackage{epsfig}
\usepackage{acronym}
\usepackage{multirow, array} 
\usepackage{float}
\usepackage[T1]{fontenc}
\usepackage{natbib}
\usepackage{fancyhdr}
\usepackage{xcolor}
\usepackage{multicol}
\usepackage{mdframed}
\usepackage{endnotes}
\usepackage{listings}
\usepackage{tensor}
\usepackage{csquotes}
\renewcommand{\thefootnote}{\arabic{footnote}}
\newcommand{\chref}[2]{
  \href{#1}{{\usebeamercolor[bg]{Feather}#2}}
}
\title[]
{
      \textbf{Título de la presentación o tarea}
}

\subtitle[NdA] %Aquí va un acrónimo de la asignatura
{Nombre de la asignatura} %Nombre de la asignatura

\author[N. Apellido] %Iniciales o nombre corto del autor
{      Nombre1 Nombre2 Apellido1 Apellido2 %Nombre completo
%Si son más de dos autores, poner \\ y luego el otro nombre para que aparezca abajo
}

\institute[]
{
      \large{Escuela Nacional de Estudios Superiores Juriquilla}\\
      Universidad Nacional Autónoma de México
  
  %there must be an empty line above this line - otherwise some unwanted space is added between the university and the country (I do not know why;( )
}
\date{}
\begin{document}


{\1% % this is the name of the PDF file for the background
\begin{frame}[plain,noframenumbering] % the plain option removes the header from the title page, noframenumbering removes the numbering of this frame only
  \titlepage % call the title page information from above
\end{frame}}
\begin{frame}
  \frametitle{Titulo de diapositiva} %Aquí va un título
  Dos ejemplos de ecuaciones alineadas:
  \begin{align*}
      a x^2 + b x + c  &= 0 \\
      f(x) &= \sum_{i=1}^{n^2-3} \frac{\int_{0}^{\infty} \left( x-r\right)^\epsilon \mathrm{d}x}{x_i^4-\sqrt{z}}
  \end{align*}
  \begin{itemize}
      \item Se puede escribir en español directamente:
      \begin{quote}
          La palabra \textbf{pedigüeñería} tiene los cuatro firuletes que un término puede tener en nuestro idioma: la virgulilla de la ñ, la diéresis sobre la ü, la tilde del acento y el punto sobre la i.
        \end{quote}
       \item ¿Acepta signos propios del español? ¡Claro que sí!
       \item No olvides cambiar la revisión de ortografía a \emph{spanish}
  \end{itemize}
\end{frame}

\begin{frame}
  \frametitle{Otro titulo de diapositiva} %Aquí va un título
\begin{enumerate}
    \item En el término \textcolor{blue}{Centrifugados}, todas las letras son diferentes y ninguna se repite, y contiene las 5 vocales.
    \item La palabra \textcolor{pink}{Oía} tiene tres sílabas en tres letras.
    \item En \textcolor{magenta}{Aristocráticos}, cada letra aparece dos veces.
    \item  El vocablo \textcolor{red}{Cinco} tiene a su vez cinco letras; la frase ``\textcolor{brown}{las doce letras}'' tiene 12 letras.
    \item Las palabras \textcolor{olive}{Ecuatorianos} y \textcolor{green}{Aeronáuticos} poseen las mismas letras, pero en diferente orden.
    \item El término \textcolor{purple}{Estuve} contiene cuatro letras consecutivas por orden alfabético: s-t-u-v.
    \item El vocablo \textcolor{orange}{Reconocer} se lee lo mismo de izquierda a derecha que viceversa (palíndromo).
    \item La palabra \textcolor{cyan}{Euforia} tiene las 5 vocales y sólo 2 consonantes.
\end{enumerate}
\begin{block}{Nota}
    Otras curiosidades del idioma español, ilustradas con colores
\end{block}
\end{frame}

\end{document}
